\documentclass[9pt,twoside]{pnas-new}
\usepackage{lastpage}
\usepackage{amsmath}
\usepackage{amssymb}
\usepackage{booktabs}
\usepackage{tabularx}
\templatetype{pnassupportinginfo}

% Clear PNAS line-counter headers: pnassupportinginfo does not suppress
% them, causing a \loop crash in the SI one-column layout.
\AtBeginDocument{\fancyhead{}}

\title{A second-order synaptic cascade motif for reconfigurable functional connectivity across nervous systems}

\author{Dongmyeong Lee, Euisun Kim, Richard Rosch, and Hae-Jeong Park}

\correspondingauthor{Hae-Jeong Park. E-mail: hjpark0@gmail.com}

\begin{document}

\maketitle

\SItext


%% ============================================================
\subsection*{Alignment of Synapse and Neuron Coordinates}
%% ============================================================

Neuronal soma coordinates were obtained from the whole-brain calcium imaging dataset \citep{chen2018brain}, in which individual neurons were localized and registered to the Z-Brain atlas. Synapse-like endpoints were obtained from the cellular-resolution anatomical reconstruction by Kunst et al.\ \citep{kunst2019cellular}, which provides the three-dimensional positions of neurite terminal points across the larval zebrafish brain.

Because the two datasets originate from independent experimental preparations, their coordinate frames were aligned using a two-step procedure. First, both datasets were registered to a common reference atlas using affine transformations, matching prominent anatomical landmarks (ventricles, olfactory bulb boundaries, and hindbrain rhombomere boundaries) visible in both datasets. Second, residual non-linear distortions were corrected by applying a thin-plate spline warp estimated from a set of manually identified point correspondences distributed across the brain volume.

%% ============================================================
\subsection*{Synaptic Classification: Pre- vs.\ Post-synaptic Endpoints}
%% ============================================================

For synapse points, terminal endpoints were classified as putative presynaptic (axonal) or postsynaptic (dendritic) based on their distance from the soma. Distal endpoints---those located far from the soma along the neurite arbor---were classified as presynaptic, because axonal boutons are found at the tips of long-range projections. Proximal endpoints---those located close to the soma---were classified as postsynaptic, reflecting the distribution of dendritic inputs.

The distance threshold separating the two classes was determined empirically by examining the bimodal distribution of soma-to-endpoint distances across all reconstructed neurons (Fig.~\ref{fig:SI_synapse}A). The threshold was set at the minimum of the density between the two modes. 

\subsection*{Sensitivity analyses support threshold-robust inter-regional structural embedding}

Sensitivity analyses confirmed that the main threshold-dependent structural result---the division-level difference in $\mathrm{DCA}_{\mathrm{post}}$---and the unadjusted regional $\mathrm{DCA}_{\mathrm{post}}$--FCV correlation were robust across a $\pm 20\,\mu$m range around the chosen threshold (Pearson $r$ between $\mathrm{DCA}_{\mathrm{post}}$ and FCV ranged from $0.580$ to $0.650$; all $p < 0.001$). As noted in the main text, this regional association was primarily carried by anatomical division differences.

To quantify the anatomical composition of the structural link set, we classified each cellular structural link as intra-regional or inter-regional using the root-area labels of its source and target neurons. A link was considered intra-regional when both neurons belonged to the same root area, and inter-regional otherwise. Across seven animals, the structural link set was dominated by inter-regional connections: inter-regional links accounted for $95.60 \pm 0.38\%$ SD of links, whereas intra-regional links accounted for $4.40 \pm 0.38\%$ SD\@. Thus, the structural link set primarily reflected connections between anatomical regions rather than within-region links.

\subsection*{Sliding-window sensitivity of FCV estimates}

To test whether FCV estimates depended strongly on the sliding-window length, we compared FCV values across alternative window settings in each species where matched outputs were available. In zebrafish, FCV was recomputed from the same high-pass filtered regional calcium traces using short, current, and long windows. The current setting used a 20-frame window with a 5-frame step; sensitivity settings used 10/3-frame and 40/10-frame window/step combinations. FCV values were highly correlated with the current setting for both the short-window condition (Pearson $r=0.957$, Spearman $\rho=0.956$, both $p<0.001$) and the long-window condition (Pearson $r=0.982$, Spearman $\rho=0.978$, both $p<0.001$; $n=66$ root-area nodes).

In \textit{Drosophila}, FCV estimates were compared across 15/5-, 20/10-, and 30/8-frame window/step settings. At the matched side-aware Ito-region scale, FCV values were strongly preserved across all window pairs (Pearson $r=0.989$--$0.998$, Spearman $\rho=0.981$--$0.997$; $n=41$ regions). The same stability was observed at the ROI scale (Pearson $r=0.986$--$0.996$, Spearman $\rho=0.986$--$0.997$; $n=464$ ROIs).

In \textit{C.~elegans}, the current 20/8-frame high-pass FCV estimates were compared with an independently computed 60/15-frame spontaneous FCV summary. Neuron-level FCV values were positively correlated across these settings (Pearson $r=0.696$, Spearman $\rho=0.685$, both $p<0.001$; $n=120$ neurons).

Together, these analyses indicate that the relative ordering of high-FCV nodes was robust to reasonable changes in sliding-window length. Thus, the main FCV patterns were not driven by a single arbitrary window choice.



\subsection*{Transfer entropy and neighborhood transfer entropy}

For each subject, neuronal calcium activity was first grouped by anatomical region. Regions containing fewer than 200 neurons were excluded. Within each retained region, we computed a neuron-by-neuron functional connectivity matrix from calcium activity traces. The absolute FC matrix was thresholded at $|r| \geq 0.2$, and Louvain community detection was applied with resolution $0.9$ to identify intra-region functional communities. Community activity was defined as the mean calcium trace of neurons assigned to the same community. Communities containing 30 or fewer neurons, or showing low temporal activity variance ($\mathrm{s.d.}<0.1$), were excluded.

Directed interactions between retained communities were quantified using transfer entropy (TE). Community mean traces were $z$-scored and discretized into four quantile-based bins. For each ordered community pair $(i,j)$, TE was computed at a lag of one frame as
\begin{equation}
\mathrm{TE}_{i \rightarrow j}
=
H(Y_{t+1},Y_t)-H(Y_t)
-
H(Y_{t+1},Y_t,X_t)+H(Y_t,X_t),
\end{equation}
where $X$ and $Y$ denote source and target community states, respectively, and $H(\cdot)$ denotes Shannon entropy in bits. Pairwise TE significance was assessed by permuting the source time series 200 times; pairs with permutation $p<0.05$ were labeled significant.

Directional asymmetry was quantified as net TE:
\begin{equation}
(\mathrm{TE}_{\mathrm{net}})_{i,j}
=
\mathrm{TE}_{i\rightarrow j}
-
\mathrm{TE}_{j\rightarrow i}.
\end{equation}
Positive values indicate stronger directed influence from community $i$ to community $j$ than in the reverse direction. For each community, outgoing net TE drive was computed as the mean net TE from that community to all other retained communities:
\begin{equation}
\mathrm{TE}_i =
\frac{1}{N-1}
\sum_{j\neq i}
(\mathrm{TE}_{\mathrm{net}})_{i,j}.
\end{equation}

To quantify local functional context, we computed neighborhood net TE drive using significant FC neighbors. Pairwise Pearson correlations between community mean traces were converted to two-sided $p$ values using the Student's $t$ distribution, and significant FC neighbors were defined after Benjamini--Hochberg FDR correction ($q<0.001$). For each community $i$, neighborhood net TE drive was defined as
\begin{equation}
\mathrm{TE}_i^{\mathrm{neigh}}
=
\frac{1}{|\mathcal{N}_i|}
\sum_{j\in\mathcal{N}_i}
\mathrm{TE}_j,
\quad
\mathcal{N}_i =
\{j:j\neq i,\ \mathrm{FC}_{i,j}\ \mathrm{significant\ after\ BH\text{-}FDR}\}.
\end{equation}
Communities with no significant FC neighbors were assigned missing values. Region-level Net TE and neighborhood Net TE values were obtained by averaging community-level values within each anatomical region and then $z$-scoring across regions.


%% ============================================================
\subsection*{Four-layer Directed Linear Network Model: Full Specification}
%% ============================================================

We considered a four-layer directed stochastic network with layer sizes
$(n_1,n_2,n_3,n_4)=(2,3,4,5)$ ($N=14$). Node activities evolved according to

\[
\frac{d\mathbf{x}}{dt}
=
J(\epsilon)\mathbf{x}(t)
+
\mathbf{u}_{\mathrm{slow}}(t)
+
0.15\,\boldsymbol{\eta}(t),
\]

where $\boldsymbol{\eta}(t)$ is Gaussian white noise. The Jacobian used the
convention that $J_{ij}$ denotes the coupling from source node $j$ to target
node $i$, with diagonal decay

\[
J_{ii}=-12.7.
\]

Within each layer, neighboring nodes were coupled symmetrically,

\[
J_{m+1,m}=J_{m,m+1}=0.25.
\]

Adjacent layers were connected all-to-all with baseline strength
$w_{\mathrm{inter}}=1.85$. Feedforward and feedback couplings were

\[
J_{ij}
=
1.85(1+s_\ell\epsilon),
\qquad
j\in L_\ell,\ i\in L_{\ell+1},
\]

\[
J_{ij}
=
1.85(1-s_\ell\epsilon),
\qquad
j\in L_{\ell+1},\ i\in L_\ell,
\]

with

\[
(s_1,s_2,s_3)=(0.8,0.6,0.1).
\]

Thus,

\[
\begin{aligned}
\mathrm{L1}\leftrightarrow\mathrm{L2}&:
(1.85+1.48\epsilon,\;1.85-1.48\epsilon),\\
\mathrm{L2}\leftrightarrow\mathrm{L3}&:
(1.85+1.11\epsilon,\;1.85-1.11\epsilon),\\
\mathrm{L3}\leftrightarrow\mathrm{L4}&:
(1.85+0.185\epsilon,\;1.85-0.185\epsilon).
\end{aligned}
\]

At $\epsilon=1$, the corresponding feedforward/feedback strengths were
$(3.33,0.37)$, $(2.96,0.74)$, and $(2.04,1.67)$.

Each layer received a slow sinusoidal input

\[
u_i(t)=\sin(2\pi\times0.1\,t+\phi_\ell),
\]

with unit amplitude and independently sampled phase $\phi_\ell$.

Dynamics were simulated using the Euler--Maruyama method
($\Delta t=0.005$, 6000 steps, 30 model time units). Functional connectivity
was estimated from sliding-window Pearson correlations of layer-averaged
activity (window = 3.0, step = 1.5). Statistics were computed for

\[
\epsilon=0,0.05,\ldots,1.0
\]

(21 values) using 30 independent simulations per value.




%% ============================================================
\subsection*{Large-scale Neuronal Simulations: Model Equations}
%% ============================================================

\subsubsection*{Spiking neuronal model}

Neuronal activity was generated using a large-scale spiking network model based on the Izhikevich formalism \citep{izhikevich2003simple}. Each neuron was described by a two-dimensional system of ordinary differential equations,
\begin{align}
\frac{dv_{i,k}}{dt} &= 0.04v_{i,k}^2 + 5v_{i,k} + 140 - u_{i,k}
+ I^{\mathrm{exc}}_{i,k} + I^{\mathrm{inh}}_{i,k} + \sigma_{i,k}, \\
\frac{du_{i,k}}{dt} &= a(bv_{i,k} - u_{i,k}),
\end{align}
with spike-reset conditions
\begin{equation}
\text{if } v_{i,k} \geq 30~\mathrm{mV}, \quad v_{i,k} \leftarrow c,\quad u_{i,k} \leftarrow u_{i,k} + d.
\end{equation}
Excitatory neurons used parameters $(a,b,c,d)=(0.02,0.2,-65,8)$; inhibitory neurons used $(0.1,0.2,-65,2)$. Gaussian noise $\sigma_{i,k}$ had zero mean and standard deviation $5$.

\subsubsection*{Synaptic currents and conduction delays}

Excitatory and inhibitory synaptic currents were given by
\begin{align}
I^{\mathrm{exc}}_{i,k} &= G_{\mathrm{exc}} g^{\mathrm{exc}}_{i,k}(10 - v_{i,k}), \\
I^{\mathrm{inh}}_{i,k} &= G_{\mathrm{inh}} g^{\mathrm{inh}}_{i,k}(-80 - v_{i,k}),
\end{align}
with $G_{\mathrm{exc}}=1$ and $G_{\mathrm{inh}}=8$. Synaptic conductances evolved according to
\begin{align}
\frac{dg^{\mathrm{exc}}_{i,k}}{dt} &= -\frac{g^{\mathrm{exc}}_{i,k}}{\tau_{\mathrm{exc}}}
+ \sum_{j \in \mathrm{exc}} w_{ij} \sum_l \delta(t - t_{j}^{l}), \\
\frac{dg^{\mathrm{inh}}_{i,k}}{dt} &= -\frac{g^{\mathrm{inh}}_{i,k}}{\tau_{\mathrm{inh}}}
+ \sum_{j \in \mathrm{inh}} w_{ij} \sum_l \delta(t - t_{j}^{l}),
\end{align}
with $\tau_{\mathrm{exc}}=1.25$ ms and $\tau_{\mathrm{inh}}=10$ ms. The spike conduction velocity was set to $2~\mathrm{m/s}$.

\subsubsection*{Calcium observation model}

Simulated membrane potentials were transformed into intracellular calcium concentrations using a phenomenological model \citep{rahmati2016inferring,yaksi2006reconstruction}:
\begin{align}
\frac{d[\mathrm{Ca}^{2+}]}{dt} &= \kappa_{\mathrm{Ca}} \, \sigma(v) - \frac{[\mathrm{Ca}^{2+}]}{\tau_{\mathrm{Ca}}}, \\
\sigma(v) &= \frac{1}{1 + e^{-R(v - V_0)}},
\end{align}
with $\kappa_{\mathrm{Ca}} = 5.8$, $\tau_{\mathrm{Ca}} = 20$, $R = 1.1$, $V_0 = 34$; parameters were held constant across simulations.

\subsubsection*{Numerical integration}

All equations were integrated with a forward Euler scheme ($\Delta t = 0.05$ ms, $T = 200$ s). Simulations were implemented in custom C/C++ code accelerated using CUDA-based GPU computation. Simulated membrane potentials were sampled at 10~Hz; functional connectivity was estimated from a 2~s sliding window with 1~s step size. Each subject-specific network was simulated five times independently, and results were averaged within animals before group-level analysis.

%% ============================================================
\subsection*{Large-scale Network Simulation: Weight Optimization and Empirical Validation}
%% ============================================================

\subsubsection*{Synaptic weight initialization}

Synaptic weights were initialized based on the empirically observed relationship between inter-neuronal Euclidean distance and functional correlation strength (Fig.~\ref{fig:SI_synapse}B). For each directed synaptic connection from neuron $j$ to neuron $i$, the initial weight was set as
\begin{equation}
  w_{ij}^{(0)} = w_0 \, \exp\!\left(-\frac{d_{ij}}{\lambda_w}\right),
\end{equation}
where $\lambda_w$ is a spatial decay constant estimated from the empirical distance--correlation curve and $w_0$ is a scaling factor. Excitatory and inhibitory connections received the same distance-dependent profile, scaled separately by $G_{\mathrm{exc}}=1$ and $G_{\mathrm{inh}}=8$.


\subsubsection*{Correlation-based weight optimization}

Synaptic weights were refined using an iterative correlation-based procedure that combined
cell-level and region-level FC constraints. For every directed structural edge from neuron $i$ to neuron $j$, the simulator computed the
empirical edge-wise calcium correlation
\begin{equation}
  C^{\mathrm{emp}}_{ij}
  =
  \mathrm{corr}
  \left(
  x^{\mathrm{emp}}_i(t),
  x^{\mathrm{emp}}_j(t)
  \right),
\end{equation}
where $x^{\mathrm{emp}}_i(t)$ and $x^{\mathrm{emp}}_j(t)$ are empirical calcium traces of the
pre- and postsynaptic neurons.

After simulating the network with the current weights, the corresponding simulated edge-wise
correlation was computed as
\begin{equation}
  C^{\mathrm{sim}}_{ij}
  =
  \mathrm{corr}
  \left(
  x^{\mathrm{sim}}_i(t),
  x^{\mathrm{sim}}_j(t)
  \right).
\end{equation}
The simulator then updated each synaptic weight according to
\begin{equation}
  \tilde{w}_{ij}^{(t+1)}
  =
  w_{ij}^{(t)}
  +
  \eta
  w_{ij}^{(t)}
  \left(
  C^{\mathrm{emp}}_{ij}
  -
  C^{\mathrm{sim}}_{ij}
  \right),
\end{equation}
where $\eta$ is the learning rate.

In addition to this edge-wise update, calcium traces were averaged within anatomical regions
to compute empirical and simulated mean FC matrices, $F^{\mathrm{emp}}$ and
$F^{\mathrm{sim}}$, using the same sliding-window correlation procedure. For each connection,
a second region-level correction was applied using the FC mismatch between the regions
containing the pre- and postsynaptic neurons:
\begin{equation}
  w_{ij}^{(t+1)}
  =
  \tilde{w}_{ij}^{(t+1)}
  +
  \xi w_{ij}^{(t)}
  \left(
  F^{\mathrm{emp}}_{r(i),r(j)}
  -
  F^{\mathrm{sim}}_{r(i),r(j)}
  \right)
  ,
\end{equation}
where $r(i)$ and $r(j)$ denote the anatomical regions of neurons $i$ and $j$, respectively,
and $\xi$ is a random scaling factor sampled independently for each update.

Weights were clipped to be non-negative after each update:
\begin{equation}
  w_{ij}^{(t+1)} \leftarrow \max \left( w_{ij}^{(t+1)}, 0 \right).
\end{equation}
To avoid excessively large total input to individual postsynaptic neurons, incoming weights
were additionally normalized when their summed input exceeded a degree-dependent upper bound.

Importantly, the optimization targeted empirical time-averaged correlation structure through
edge-wise calcium correlations and region-level mean FC\@. No explicit constraint was imposed
on temporal FC fluctuations; FCV was therefore evaluated from the fitted directed structural
topology and network dynamics.

\subsubsection*{Empirical versus simulated FC and FCV}

After optimization, simulated mean FC matrices were compared with empirical FC matrices region-by-region across all subjects (Fig.~\ref{fig:SI_sim}). Regional mean FC showed strong correspondence between empirical and simulated data ($r = 0.936$, $p < 0.001$; Fig.~\ref{fig:SI_sim}G), with lower mean FC in telencephalic regions (pallium: $-0.50 \pm 0.59$; subpallium: $-0.68 \pm 0.55$) relative to the cerebellum ($0.43 \pm 0.47$). FC variability also showed significant correspondence ($r = 0.445$, $p < 0.01$; Fig.~\ref{fig:SI_sim}H), with elevated FCV in telencephalic regions (pallium: $0.44 \pm 0.79$; subpallium: $0.67 \pm 0.85$) compared with the cerebellum ($0.20 \pm 0.71$). These results indicate that the synapse-informed directed structural network reproduces key features of large-scale functional organization and is associated with regional differences in FC variability.

\begin{table*}[!ht]
\centering
\scriptsize
\caption{
\textbf{Anatomical and functional group assignments across species.}
Anatomical groups were defined from the SC--FC analysis table.}
\label{tab:species_group_assignments}
\begin{tabularx}{\linewidth}{lllX}
\toprule
Grouping & Species & Group & Included nodes or regions \\
\midrule
\multicolumn{4}{l}{\textbf{Anatomical groups}} \\
\midrule
Anatomical & \textit{C.~elegans} & sensory neurons & ADEL, ADER, ADLL, ADLR, AQR, ASEL, ASER, ASGL, ASGR, ASHL, ASIL, ASIR, ASKL, ASKR, AWAL, AWAR, AWBL, AWBR, AWCL, AWCR, BAGL, CEPDL, CEPDR, CEPVL, CEPVR, FLPL, FLPR, IL1DL, IL1DR, IL1L, IL1R, IL1VL, IL1VR, IL2DL, IL2DR, IL2L, IL2R, IL2VL, IL2VR, OLLL, OLLR, OLQDL, OLQDR, OLQVL, OLQVR, URXL, URXR, URYDL, URYDR, URYVL, URYVR \\
Anatomical & \textit{C.~elegans} & interneurons & ADAL, ADAR, AIBL, AIBR, AIML, AIMR, AINL, AINR, AIYL, AIYR, AIZL, AIZR, AUAL, AUAR, AVHL, AVJL, AVJR, AVKL, AVKR, AVL, RIAL, RIAR, RIBL, RIBR, RID, RIH, RIR, URBL, URBR \\
Anatomical & \textit{C.~elegans} & motor / command & AVAL, AVAR, AVBL, AVBR, AVDL, AVDR, AVEL, AVER, RIVL, RIVR, RMDDL, RMDDR, RMDL, RMDR, RMDVL, RMDVR, RMED, RMEL, RMER, RMEV, SAADL, SAADR, SAAVL, SAAVR, SMBVR, SMDDL, SMDDR, SMDVL, SMDVR, URADL, URADR, URAVL, URAVR \\
Anatomical & \textit{C.~elegans} & state-modulatory & ALA, RICL, RICR, RIML, RIMR, RIS, RMGL \\

Anatomical & \textit{Drosophila} & olfactory neuropil & AL\_R, LH\_L, LH\_R \\
Anatomical & \textit{Drosophila} & optic neuropil & AME\_L, LO\_R, ME\_L, ME\_R \\
Anatomical & \textit{Drosophila} & mushroom body & MB\_L, MB\_R \\
Anatomical & \textit{Drosophila} & premotor interface & LAL\_L, LAL\_R, SPS\_L, SPS\_R, VES\_L, VES\_R \\
Anatomical & \textit{Drosophila} & protocerebrum & ATL\_L, ATL\_R, AVLP\_L, AVLP\_R, CRE\_L, CRE\_R, EPA\_L, EPA\_R, IB\_L, IB\_R, ICL\_L, ICL\_R, PLP\_L, PLP\_R, PVLP\_L, PVLP\_R, SIP\_L, SIP\_R, SLP\_R, SMP\_L, SMP\_R, WED\_R \\

Anatomical & zebrafish & telencephalon & OB, P, SP, rOB, rP, rSP \\
Anatomical & zebrafish & diencephalon & PT, PrT, T, Th, rPT, rPrT, rT, rTh \\
Anatomical & zebrafish & mesencephalon & TL, TS, TeO, rTS, rTeO \\
Anatomical & zebrafish & hindbrain / motor & Cb, MON, MOS1, MOS2, MOS3, MOS4, MOS5, NX, aRF, imRF, pRF, rCb, rMON, rMOS1, rMOS2, rMOS3, rMOS4, rMOS5, rNX, raRF, rimRF, rpRF \\
Anatomical & zebrafish & state system & Hb, PO, Ra, rHb, rPO, rRa \\

\midrule
\multicolumn{4}{l}{\textbf{Functional groups}} \\
\midrule
Functional & \textit{C.~elegans} & olf./chemo & ADLL, ADLR, ASEL, ASER, ASGL, ASGR, ASHL, ASIL, ASIR, AWAL, AWAR, AWBL, AWBR, AWCL, AWCR \\
Functional & \textit{C.~elegans} & other sensory & ADEL, ADER, AQR, ASKL, ASKR, BAGL, CEPDL, CEPDR, CEPVL, CEPVR, FLPL, FLPR, IL1DL, IL1DR, IL1L, IL1R, IL1VL, IL1VR, IL2DL, IL2DR, IL2L, IL2R, IL2VL, IL2VR, OLLL, OLLR, OLQDL, OLQDR, OLQVL, OLQVR, URXL, URXR, URYDL, URYDR, URYVL, URYVR \\
Functional & \textit{C.~elegans} & sensorimotor & AVAL, AVAR, AVBL, AVBR, AVDL, AVDR, AVEL, AVER, RIVL, RIVR, RMDDL, RMDDR, RMDL, RMDR, RMDVL, RMDVR, RMED, RMEL, RMER, RMEV, SAADL, SAADR, SAAVL, SAAVR, SMBVR, SMDDL, SMDDR, SMDVL, SMDVR, URADL, URADR, URAVL, URAVR \\
Functional & \textit{C.~elegans} & integrative & ADAL, ADAR, AVHL, AVJL, AVJR, AVKL, AVKR, AVL, RID, RIH, RIR, URBL, URBR \\
Functional & \textit{C.~elegans} & assoc./state & AIBL, AIBR, AIML, AIMR, AINL, AINR, AIYL, AIYR, AIZL, AIZR, ALA, AUAL, AUAR, RIAL, RIAR, RIBL, RIBR, RICL, RICR, RIML, RIMR, RIS, RMGL \\

Functional & \textit{Drosophila} & olf./chemo & AL\_L, AL\_R, LH\_L, LH\_R \\
Functional & \textit{Drosophila} & visual & AME\_L, AME\_R, LO\_L, LO\_R, ME\_L, ME\_R \\
Functional & \textit{Drosophila} & sensorimotor & LAL\_L, LAL\_R, SPS\_L, SPS\_R, VES\_L, VES\_R \\
Functional & \textit{Drosophila} & integrative & AMMC\_L, AMMC\_R, ATL\_L, ATL\_R, AVLP\_L, AVLP\_R, CRE\_L, CRE\_R, EPA\_L, EPA\_R, FLA, GA\_L, GA\_R, GOR\_L, GOR\_R, IB\_L, IB\_R, ICL\_L, ICL\_R, IPS\_L, IPS\_R, IVLP\_L, IVLP\_R, LB\_L, LB\_R, LOP\_R, OTU\_L, OTU\_R, PLP\_L, PLP\_R, PRW, PVLP\_L, PVLP\_R, SAD, SCL\_L, SCL\_R, SIP\_L, SIP\_R, SLP\_L, SLP\_R, SMP\_L, SMP\_R, SOG, WED\_L, WED\_R \\
Functional & \textit{Drosophila} & assoc./state & EB, FB, MB\_L, MB\_R, NO, PB \\

Functional & zebrafish & olf./chemo & OB \\
Functional & zebrafish & visual & PT, PrT, TS, TeO, rPT, rPrT, rTS, rTeO \\
Functional & zebrafish & sensorimotor & Cb, MON, MOS1, MOS2, MOS3, MOS4, MOS5, NX, aRF, imRF, pRF, rCb, rMON, rMOS1, rMOS2, rMOS3, rMOS4, rMOS5, rNX, raRF, rimRF, rpRF \\
Functional & zebrafish & integrative & T, Th, rT, rTh \\
Functional & zebrafish & assoc./state & HR, Hb, Hi, P, PO, Ra, SP, rHR, rHb, rHi, rP, rPO, rRa, rSP \\
\bottomrule
\end{tabularx}
\end{table*}




\begin{table}[h]
\centering
\small
\begin{tabular}{lll}
\hline
Major division & Abbreviation & Brain region name \\ \hline
Telencephalon & L\_P, R\_P & Pallium \\
Telencephalon & L\_SP, R\_SP & Subpallium \\
Telencephalon & L\_OB, R\_OB & Olfactory Bulb \\
Telencephalon & L\_PO, R\_PO & Preoptic Region \\
Diencephalon & L\_Hb, R\_Hb & Habenula \\
Diencephalon & L\_Hc, R\_Hc & Caudal Hypothalamus \\
Diencephalon & L\_Hi, R\_Hi & Intermediate Hypothalamus \\
Diencephalon & L\_Hr, R\_Hr & Rostral Hypothalamus \\
Diencephalon & L\_PT, R\_PT & Posterior Tuberculum \\
Diencephalon & L\_Th, R\_Th & Thalamus \\
Diencephalon & L\_Pi, R\_Pi & Pituitary \\
Mesencephalon & L\_PrT, R\_PrT & Pretectum \\
Mesencephalon & L\_TeO, R\_TeO & Tectum \\
Mesencephalon & L\_T, R\_T & Tegmentum \\
Mesencephalon & L\_TL, R\_TL & Torus Longitudinalis \\
Mesencephalon & L\_TS, R\_TS & Torus Semicircularis \\
Hindbrain & L\_Cb, R\_Cb & Cerebellum \\
Hindbrain & L\_IPN, R\_IPN & Interpeduncular Nucleus \\
Hindbrain & L\_IO, R\_IO & Inferior Olive \\
Hindbrain & L\_MON, R\_MON & Medial Octavolateral Nucleus \\
Hindbrain & L\_MOS1--5, R\_MOS1--5 & Medulla Oblongata Strips 1--5 \\
Hindbrain & L\_Ra, R\_Ra & Raphe Nucleus \\
Hindbrain & L\_aRF, R\_aRF & Anterior Reticular Formation \\
Hindbrain & L\_imRF, R\_imRF & Intermediate Reticular Formation \\
Hindbrain & L\_pRF, R\_pRF & Posterior Reticular Formation \\
Hindbrain & L\_VR, R\_VR & Vagal Region \\
Hindbrain & L\_NX, R\_NX & Vagus Motor Neurons \\ \hline
\end{tabular}
\caption{\textbf{Brain regions included in the analysis.} Complete list with abbreviations and major divisions following established zebrafish atlases \citep{Randlett2015ZBrain,Wullimann1996Neuroanatomy,mueller2015atlas}. L\_ and R\_ prefixes denote left and right hemispheres.}
\label{tab:regions}
\end{table}


\begin{table}[h]
\centering
\small
\begin{tabular}{lll}
\hline
Major division & Abbreviation & Brain region name \\ \hline
Olfactory neuropil & AL\_L, AL\_R & Antennal Lobe \\
Olfactory neuropil & LH\_L, LH\_R & Lateral Horn \\

Optic neuropil & AME\_L, AME\_R & Accessory Medulla \\
Optic neuropil & LO\_L, LO\_R & Lobula \\
Optic neuropil & ME\_L, ME\_R & Medulla \\

Associative / state-related & MB\_L, MB\_R & Mushroom Body \\
Associative / state-related & EB & Ellipsoid Body \\
Associative / state-related & FB & Fan-shaped Body \\
Associative / state-related & NO & Noduli \\
Associative / state-related & PB & Protocerebral Bridge \\

Premotor interface & LAL\_L, LAL\_R & Lateral Accessory Lobe \\
Premotor interface & SPS\_L, SPS\_R & Superior Posterior Slope \\
Premotor interface & VES\_L, VES\_R & Vest \\

Protocerebrum / integrative neuropil & ATL\_L, ATL\_R & Antler \\
Protocerebrum / integrative neuropil & AVLP\_L, AVLP\_R & Anterior Ventrolateral Protocerebrum \\
Protocerebrum / integrative neuropil & CRE\_L, CRE\_R & Crepine \\
Protocerebrum / integrative neuropil & EPA\_L, EPA\_R & Epaulette \\
Protocerebrum / integrative neuropil & IB\_L, IB\_R & Inferior Bridge \\
Protocerebrum / integrative neuropil & ICL\_L, ICL\_R & Inferior Clamp \\
Protocerebrum / integrative neuropil & PLP\_L, PLP\_R & Posterior Lateral Protocerebrum \\
Protocerebrum / integrative neuropil & PVLP\_L, PVLP\_R & Posterior Ventrolateral Protocerebrum \\
Protocerebrum / integrative neuropil & SIP\_L, SIP\_R & Superior Intermediate Protocerebrum \\
Protocerebrum / integrative neuropil & SLP\_L, SLP\_R & Superior Lateral Protocerebrum \\
Protocerebrum / integrative neuropil & SMP\_L, SMP\_R & Superior Medial Protocerebrum \\
Protocerebrum / integrative neuropil & WED\_L, WED\_R & Wedge \\

Other integrative neuropil & AMMC\_L, AMMC\_R & Antennal Mechanosensory and Motor Center \\
Other integrative neuropil & FLA & Flange \\
Other integrative neuropil & GA\_L, GA\_R & Gall \\
Other integrative neuropil & GOR\_L, GOR\_R & Gorget \\
Other integrative neuropil & IPS\_L, IPS\_R & Inferior Posterior Slope \\
Other integrative neuropil & IVLP\_L, IVLP\_R & Inferior Ventrolateral Protocerebrum \\
Other integrative neuropil & LB\_L, LB\_R & Bulb \\
Other integrative neuropil & LOP\_R & Lobula Plate \\
Other integrative neuropil & OTU\_L, OTU\_R & Optic Tubercle \\
Other integrative neuropil & PRW & Prow \\
Other integrative neuropil & SAD & Saddle \\
Other integrative neuropil & SCL\_L, SCL\_R & Superior Clamp \\
Other integrative neuropil & SOG & Subesophageal Ganglion \\ \hline
\end{tabular}
\caption{\textbf{\textit{Drosophila} brain regions included in the analysis.}
Side-aware Ito-region abbreviations used for matching FlyWire structural connectivity and Branson calcium imaging data. L\_ and R\_ prefixes denote left and right hemisphere labels when available. Region names follow the adult fly brain nomenclature \citep{Ito2014InsectBrainNomenclature,Dorkenwald2024FlyWire,Turner2021Branson999}.}
\label{tab:drosophila_regions}
\end{table}



%% ============================================================
%%  Supporting Figures
%% ============================================================

\begin{figure}[p]
\centering
\includegraphics[width=\linewidth]{supporting_figures/figure_supply_1.png}
\caption{
\textbf{Synapse classification and connectome construction quality control.}
(A) Distribution of soma-to-endpoint distances across all reconstructed neurons, used to classify terminal endpoints as putative presynaptic (distal, red) or postsynaptic (proximal, blue). The dashed vertical line indicates the classification threshold; the shaded region denotes the $\pm 10\,\mu$m exclusion zone.
(B) Empirical relationship between inter-neuronal Euclidean distance and functional correlation strength, pooled across all subjects ($n=7$). The solid curve shows the exponential fit used to estimate $\lambda = 80\,\mu$m for the probabilistic connection model and to initialize synaptic weights. Shaded region denotes $95\%$ confidence interval.
}
\label{fig:SI_synapse}
\end{figure}

\begin{figure}[p]
\centering
\includegraphics[width=\linewidth]{supporting_figures/figure_supply_2_proc.png}
\caption{
\textbf{Regional structural embedding measures across major brain divisions.}
(A) $\mathrm{DCA}_{\mathrm{post}}$, (B) $\mathrm{DCA}_{\mathrm{pre}}$, (C) modularity ($Q$), (D) log$_{10}$(out/in degree ratio), and (E) output--output (OO) motif fraction for anatomically defined zebrafish brain regions.
Bars represent mean values across animals and error bars indicate SEM.
Region labels are ordered to match Fig.~\ref{fig:SI_fc_metrics} and colored by major brain division: telencephalon (green), diencephalon (orange), mesencephalon (red), and hindbrain (blue).
Telencephalic regions are highlighted in bold italic.
These measurements summarize regional variation in directional downstream and upstream embedding, mesoscale modular organization, first-order output/input imbalance, and second-order output-to-output cascade motifs in the structural connectome.
}
\label{fig:SI_struct_metrics}
\end{figure}


\begin{figure}[p]
\centering
\includegraphics[width=\linewidth]{supporting_figures/figure_supply_10_proc.png}
\caption{
\textbf{Regional functional dynamics and directed information flow across major brain divisions.}
(A) Functional connectivity variability (FCV), (B) functional connectivity strength (FCS), (C) FC partner reconfiguration, (D) net transfer entropy ($\mathrm{TE}_{\mathrm{net}}$), and (E) neighbor-restricted $\mathrm{TE}_{\mathrm{net}}$ for anatomically defined zebrafish brain regions.
Bars represent mean values across animals and error bars indicate SEM.
Region labels are colored by major brain division: telencephalon (green), diencephalon (orange), mesencephalon (red), and hindbrain (blue).
Telencephalic regions are highlighted in bold italic.
These measurements summarize regional differences in temporal FC variability, mean coupling strength, FC partner reconfiguration, directed information flow, and directed drive among functionally connected neighbors.
}
\label{fig:SI_fc_metrics}
\end{figure}


%\begin{figure}[p]
%\centering
%\includegraphics[width=\linewidth]{supporting_figures/figure_supply_11.png}
%\caption{
%\textbf{Robustness and division-control analyses for structural prediction of regional FCV.}
%(A) Permutation-null test for the linear cross-validated prediction of FCV from structural-connectivity features. The observed 5-fold cross-validated performance was $R^2 = 0.334$ ($r = 0.595$, $p = 1.66 \times 10^{-5}$), exceeding the null distribution generated by permuting FCV values across regions ($n = 10{,}000$ permutations; null mean $R^2 = -0.283$, 95th percentile $R^2 = -0.018$, permutation $p = 2.0 \times 10^{-4}$).
%(B) Drop-one-feature analysis for the linear model. Each structural feature was removed one at a time, and the decrease in cross-validated $R^2$ was measured relative to the full model. Removing $\mathrm{DCA}_{\mathrm{post}}$ caused the largest loss of performance ($\Delta R^2 = 0.187$), whereas removal of other features produced much smaller or negative changes ($\Delta R^2$ range: $-0.020$ to $0.012$).
%(C) Residual proxy analysis testing whether absolute residuals of the linear SC-to-FCV model were associated with measurement or sampling proxies. Residual magnitude was not significantly associated with FCV SEM ($r = 0.132$, $p = 0.389$), calcium trace variance SEM ($r = -0.085$, $p = 0.581$), or mean neuron count ($r = -0.129$, $p = 0.399$).
%(D) Subject-level replication of the $\mathrm{DCA}_{\mathrm{post}}$--FCV association. Pearson correlations were computed separately for each animal using the same region-level $\mathrm{DCA}_{\mathrm{post}}$ values and animal-specific FCV estimates. The random-effects Fisher-$z$ summary showed a consistent positive association across animals ($r = 0.473$, 95\% CI $= 0.380$--$0.557$, $p = 1.12 \times 10^{-18}$; $\tau^2 = 0$).
%(E) Leave-one-animal-out prediction. For each held-out animal, the model was fit using the mean FCV pattern from the remaining animals and then used to predict FCV in the held-out animal. Across all held-out region-animal observations, predicted and observed FCV were correlated ($r = 0.473$, $R^2 = 0.221$, $p = 5.55 \times 10^{-19}$).
%(F) After residualizing both $\mathrm{DCA}_{\mathrm{post}}$ and FCV with respect to anatomical division, the residual-residual association was weak and not significant ($r_{\mathrm{partial}} = 0.123$, $p = 0.420$). Each point represents one brain region, colored by anatomical division.
%(G) A null distribution was generated by shuffling $\mathrm{DCA}_{\mathrm{post}}$ values within each anatomical division ($n = 10{,}000$ permutations). The observed division-controlled association did not exceed the within-division permutation null ($p_{\mathrm{perm}} = 0.486$).
%(H) A division-only model explained most of the FCV variance ($R^2 = 0.682$). Adding standardized $\mathrm{DCA}_{\mathrm{post}}$ produced only a small increase in explained variance ($R^2 = 0.686$, $\Delta R^2 = 0.00484$; $\mathrm{DCA}_{\mathrm{post}}$ $\beta = 0.0947$).
%Structural features predict regional FCV above permutation-null expectations and are consistent across animals, with $\mathrm{DCA}_{\mathrm{post}}$ contributing strongly to the original multivariate prediction. Division-control analyses indicate that the $\mathrm{DCA}_{\mathrm{post}}$--FCV association is explained primarily by broad anatomical division structure rather than by a strong continuous within-division relationship.
%}
%\label{fig:SI_fc_vs_sc_metrics}
%\end{figure}


\begin{figure}[p]
\centering
\includegraphics[width=\linewidth]{supporting_figures/figure_supply_13.png}
\caption{
\textbf{Representative subject-level directed functional interaction estimated by transfer entropy.}
(A) Net transfer entropy ($\mathrm{TE}_{\mathrm{net}}$) matrix from subject 13, sorted by major anatomical division. $\mathrm{TE}_{\mathrm{net}}$ was calculated as the directional difference between pairwise transfer entropy values, $(\mathrm{TE}_{\mathrm{net}})_{i \rightarrow j} = \mathrm{TE}_{i \rightarrow j} - \mathrm{TE}_{j \rightarrow i}$, so positive values indicate stronger directed influence from source to target. Black lines separate telencephalon, diencephalon, mesencephalon, and hindbrain communities.
(B) Example activity traces for a representative high-Net-TE pair originating from the P region. The displayed pair corresponds to P community 9 projecting to MOS1 community 61; traces are shown as z-scored calcium activity.
(C) Permutation-null test for the forward TE of the displayed pair. The observed forward TE was 0.0678 bits, exceeding the null distribution generated by permuting the source activity trace ($n=1{,}000$ permutations; null mean = 0.0230 bits, 95th percentile = 0.0321 bits, permutation $p = 0.001$).
(D) Directional comparison of transfer entropy for the same pair. Forward TE from P to MOS1 was larger than the reverse direction (forward TE = 0.0678 bits; reverse TE = 0.0128 bits), yielding a positive $\mathrm{TE}_{\mathrm{net}}$ of 0.0550 bits.
This example illustrates how asymmetric temporal dependencies in regional calcium activity were used to estimate directed functional interactions, showing a representative P-origin interaction with stronger forward than reverse information transfer.
}
\label{fig:SI_TE}
\end{figure}



\begin{figure}[p]
\centering
\includegraphics[width=\linewidth]{supporting_figures/figure_supply_14.png}
\caption{
\textbf{Stimulus-evoked FCV profiles across brain regions during OMR-related conditions.}
Heatmap showing mean stimulus-evoked functional connectivity variability (FCV) for each brain region across three optomotor response (OMR)-related stimulus conditions: stimulus 10, OMR forward; stimulus 11, OMR rightward; and stimulus 12, OMR leftward. For each stimulus and region, FCV was averaged across animals. Columns indicate anatomically defined brain regions, ordered by major brain division, and rows indicate stimulus conditions. The color bar shows mean stimulus FCV\@. The top annotation bar and region-label colors indicate major brain divisions: telencephalon, diencephalon, mesencephalon, and hindbrain.
}
\label{fig:SI_stimulus_fcv}
\end{figure}



\begin{figure}[p]
\centering
\includegraphics[width=\linewidth]{supporting_figures/figure_supply_0.png}
\caption{
\textbf{Empirical and simulated functional connectivity.}
(A) Representative simulated calcium activity traces from anatomically ordered regions. Region labels are colored by major brain division.
(B) Representative time-resolved FC traces for a telencephalic connection (pallium--cerebellum, P--Cb; blue) and a hindbrain connection (cerebellum--medulla oblongata strip 2, Cb--MOS2; orange), illustrating the contrast in temporal variability.
(C) Empirical and simulated mean FC matrices for a representative subject, ordered by anatomical region.
(D) Empirical and simulated FCV matrices for the same representative subject and ordering.
(E) Region-wise empirical and simulated FCS profiles, shown as z-scored values across anatomical regions. Bars show empirical values and black markers show simulated values.
(F) Region-wise empirical and simulated FCV profiles, shown as z-scored values across anatomical regions. Bars show empirical values and black markers show simulated values.
(G) Group-level correspondence between empirical and simulated mean regional FCS, averaged across subjects (Pearson $r=0.936$, $p<0.001$). Each point represents one region; color indicates major brain division.
(H) Group-level correspondence between empirical and simulated regional FCV, averaged across subjects (Pearson $r=0.445$, $p<0.01$). Color scheme as in (G).
}
\label{fig:SI_sim}
\end{figure}




\begin{figure}[p]
\centering
\includegraphics[width=\linewidth]{supporting_figures/figure_supply_5.png}
\caption{
\textbf{Directional coreness asymmetry is selectively disrupted by null rewiring.}
(A) $\mathrm{DCA}_{\mathrm{post}}$ in the Base network and under the NULL-Out condition. NULL-Out significantly altered $\mathrm{DCA}_{\mathrm{post}}$ relative to Base ($p < 0.001$).
(B) $\mathrm{DCA}_{\mathrm{pre}}$ in the Base network and under the NULL-In condition. NULL-In significantly altered $\mathrm{DCA}_{\mathrm{pre}}$ relative to Base ($p < 0.001$).
(C) Representative network diagrams for the rP--rimRF connection under the original structural network, NULL-In rewiring, and NULL-Out rewiring, respectively.
(D) Representative network diagrams for the rP--rMOS4 connection under the original structural network, NULL-In rewiring, and NULL-Out rewiring, respectively.
Directional null rewiring selectively disrupts the directional embedding present in the baseline structural connectivity: NULL-Out preferentially affects post-synaptic directional coreness and NULL-In preferentially affects pre-synaptic directional coreness.
}
\label{fig:SI_null_dca}
\end{figure}


\begin{figure}[p]
\centering
\includegraphics[width=\linewidth]{supporting_figures/figure_supply_15.png}
\caption{
\textbf{Representative activity traces under null rewiring and asymmetric layer-linear simulation.}
(A) Representative simulated calcium activity traces for the rSP--rMOS5 pair under the Base, NULL-In, and NULL-Out structural-network conditions. The rSP trace and rMOS5 trace are shown as smoothed, z-scored activity over the same selected time window. NULL-In and NULL-Out conditions were generated by selectively rewiring incoming or outgoing connections of the P/SP-related network while preserving the corresponding null-control structure.
(B) Representative traces from the four-layer linear model under increasing between-layer directional asymmetry. Mean activity from layer 1 (L1) and layer 4 (L4) is shown for $\epsilon = 1.0$, $0.3$ and $0$. The same model parameters as in the layer-linear analysis were used, with $\epsilon$ controlling the strength of feedforward versus feedback coupling between adjacent layers.
These examples illustrate how directional perturbation of empirical structural connectivity and increasing asymmetry in a simplified layered model alter the temporal relationship between upstream and downstream activity.
}
\label{fig:SI_traces}
\end{figure}








\bibliography{reference}

\end{document}
